\documentclass[aspectratio=169,notheorems]{corvinusmetropolis}

\usepackage[hungarian]{babel}
\usepackage[style=apa, backend=biber]{biblatex}
\addbibresource{05-05.bib}
\urlstyle{tt}
\usefonttheme{serif}
\setbeamertemplate{frametitle continuation}{}

\renewcommand{\myinstlogo}{./corvinus-logos/black_simple_HUN.pdf}

\title{Biodiversity \& Nature Risk in Europe}
\subtitle{Biodiverzitás és Nature kockázat megjelenése az európai piacon, elmélet és piaci realitás}
\author{Molnár Marcell,\\Konzulens: Dömötör Barbara}
\email{marcell.molnar@stud.uni-corvinus.hu}

\institute[Corvinus University]
  {
  Pénzügy és számvitel alapszak\\
  Pénzügy specializáció\\
  Email: \putemail\\
  }
\date{\today}
\titlegraphic{\hfill{\includegraphics[width=.2\textwidth]{\myinstlogo}}}
\keywords{biodiversity_finance, nature_finance, sustainable_finance}

\begin{document}
\openingframes{}

\begin{frame}
  \frametitle{Irodalomi háttér}
  Az elmúlt években a Fenntartható Pénzügyek kutatói fokozott figyelmet fordított az új publikációkban a Biodiversity Finance ágnak\\ 
  Friss, releváns kutatások:
  \begin{itemize}
    \item \alert{biodiverzitási prémium} jelenléte és mértéke \parencite{coqueret_biodiversity_2024, naffa_biodiversity_2024}
    \item \alert{ENCORE adatbázis} biodiverzitás témakörben \parencite{clement_mapping_2025}
    \item \alert{private capital befektetések} \parencite{flammer_biodiversity_2025}\\
  \end{itemize}
  \textbf{A hozzájárulásunk a szakirodalomhoz a biodiversity/nature hatások empirikus vizsgálata az európai részvénypiacon}
\end{frame}

\begin{frame}
  \frametitle{Módszertan I.}
  Egy olyan befektetőt feltételezünk (shortolás megengedett), akinek a hasznossági függvényébe a vagyon maximalizálása mellett, \alert{pozitívan gondolkodik a magasabb biodiversity/nature értékekről}.\\
  \textcite{pedersen_responsible_2021} alapján a \alert{hasznossági fv-e lineáris és a következő alakú}:
  \metroset{block=fill}
  \begin{block}{Kiegészített hasznossági függvény}
  \begin{equation}
    U = E[\widehat{W}|s] - \frac{\hat{\gamma}}{2}Var(\widehat{W}|s) + Wf(\bar{s}),
  \end{equation}
  ahol $\widehat{W}=W(1+r^f+x'r)$, $\hat{\gamma}$ az abszolút kockázatelutasító paraméter, $s$ és $\bar{s}$ pedig a biodiversity/nature score és számtani átlaga.
  \end{block}
\end{frame}

\begin{frame}
  \frametitle{Módszertan II.}
  Egy sztenderd mean-variance hasznosságfüggvényű befektető a befektetési univerzumban megkeresi az \alert{érintési portfóliót}, majd ezt keveri a \alert{kockázatmentes eszközzel} az egyéni kockázatkerülésének függvényében.
  \\Az érintési portfólió az a portfólió, ami a legmagasabb Sharpe-rátával, ezért \textcite{pedersen_responsible_2021} \textbf{általánosította} ezt azzal, hogy \alert{minden egyes ESG-proxy mellett} (jelenesteben biodiversity score), \alert{meghatározható a maximálisan elérhető Sharpe-ráta}.\\
  \bigskip
  \textbf{A probléma idáig nem függ az egyéni hasznosságtól, azaz a befektetési univerzumra egyértelműen meghatározható az ún. BIODIV-SR határ!}
\end{frame}

\begin{frame}
  \frametitle{Módszertan III.}
  A biodiverzitási prémiumot \textcite{naffa_biodiversity_2024} után határozzuk meg:
  \metroset{block=fill}
  \begin{block}{Biodiverzitási prémium (adott $\alpha$ mellett)}
  \begin{equation}
    \text{BRP}^\alpha=\text{SR}^{\alpha}_{\text{bio}} - \overline{\text{SR}^{\alpha, k}_{\text{rnd}}}
  \end{equation}
  Itt $\text{SR}^{\alpha}_{\text{bio}}$ a BIODIV-SR adott $\alpha$ biodiversity score melletti maximum elérhető Sharpe-ráta,
  $\overline{\text{SR}^{\alpha, k}_{\text{rnd}}}$ pedig $k$-szor elvégzett visszatevés nélküli mintavéttellel készített portfóliók maximum Sharpe-rátáinak átlaga.
  \end{block}
  \alert{A random generált portfólió \textbf{eszközeinek számossága mindig meg kell egyeznie} a BIODIV-SR adott pontjaihoz tartozó portfólió számosságával}
\end{frame}

\begin{frame}
  \frametitle{Kutatási kérdések}
  \begin{enumerate}
    \item A biodiverzitási adatok beépítésével növelni tudjuk a maximálisan elérhető Sharpe-rátát, az egyszerű kiszűrési stratégiához képest? \footcite{pedersen_responsible_2021} 
    \item A piaci hozamokban megjelenik-e a biodiverzitási prémium \footcite{coqueret_biodiversity_2024, naffa_biodiversity_2024}, és ha igen, akkor ennek a méréke mekkora az európai piacon?
  \end{enumerate}
\end{frame}

\begin{frame}
  \frametitle{Adatbázisok, befektetési univerzum}
\end{frame}

\begin{frame}
  \frametitle{Roadmap}
\end{frame}

\begin{frame}[allowframebreaks]
  \frametitle{Hivatkozások}
  \printbibliography[heading=none]
\end{frame}

\closingframes{Köszönöm a figyelmet!}
\end{document}
